\section{Uniform Pad\'e approximation, Stieltjes structure, and Borel reconstruction}
\label{app:pade-borel}

This appendix collects the approximation guarantees discussed after the
population-flow theorem. They answer a separate question: when can finite
initialization data approximate an observable uniformly on a fixed interval,
even if its Taylor partial sums fail? These are sufficient theorems, not
properties already established for the neural kernel.

In the general examples, $F$ denotes a scalar target and $a_p$ its formal
coefficient of order $p$. Thus $a_p=F^{(p)}(0)/p!$ when actual derivatives
exist. The unindexed $\sigma>0$ is the Borel exponent, not the layer
initialization scales $\sigma_\ell$. The letter $R$ below bounds the
support of a measure; it is not the backward-field cutoff of Appendix A.
All approximation symbols are local to this appendix.

\subsection{Which rational approximation is being requested?}

Write
\[
 F(t)\sim\sum_{p=0}^{\infty}a_p t^p.
\]
A normalized Pad\'e approximant of type $[N/M]$ is $P_N/Q_M$, with
polynomial degrees at most $N,M$, $Q_M(0)=1$, and
\[
 Q_M(t)\sum_{p=0}^{\infty}a_p t^p-P_N(t)=O(t^{N+M+1}).
\]
This is formal coefficient agreement through the specified order, not an
assumption of series convergence. Only $a_0,\ldots,a_{N+M}$ enter.
Some Pad\'e-table entries are defective: the normalized matching problem
need not be solvable for every pair of degrees. The positive-measure
constructions below provide well-defined entries, with finite-support
cases terminating in an exact rational function.

An arbitrary continuous function on $[0,T]$ admits uniform polynomial
approximation, hence rational approximation. Those polynomials generally
use global function information. Direct Pad\'e uses only initialization
coefficients and the displayed matching equations. Borel--Pad\'e instead
applies Pad\'e to a transformed series and then uses a reconstruction
integral. It generally does not produce the direct Pad\'e approximant of $F$.

Failure of Taylor approximation here means failure of the function's own
Taylor partial sums, not failure of every polynomial approximation.
A nonzero smooth function flat at zero explains why initialization data
require additional structure: all coefficients vanish, so every normalized
Pad\'e approximant of that jet is zero. Appendix
\ref{app:borel-population} gives the identification obstruction even for a
unique smooth evolution.

\subsection{A genuine Pad\'e theorem with an explicit uniform error}
\label{app:compact-stieltjes}

\begin{proposition}[Compactly supported positive Stieltjes integral]
Suppose the actual function is
\[
 F(t)=\int_{[0,R]}\frac{d\mu(x)}{1+tx},
 \qquad \mu\ge0,\quad R>0,\quad m_0=\mu([0,R])<\infty.
\]
Its coefficients are $a_p=(-1)^pm_p$, where $m_p=\int x^p\,d\mu(x)$.
If the support is infinite, its $[M-1/M]$ Pad\'e approximant exists
for every $M\ge1$ and has the form
\[
 R_M(t)=\sum_{j=1}^M\frac{w_j}{1+tx_j},
 \qquad w_j>0,\quad 0<x_j<R,\quad \sum_jw_j=m_0.
\]
It uses only the first $2M$ coefficients, has no nonnegative real pole,
and for every finite $T>0$ satisfies
\begin{equation}
 \sup_{0\le t\le T}|F(t)-R_M(t)|
 \le 2m_0\left(\frac{RT}{2+RT}\right)^{2M}.
 \label{eq:stieltjes-uniform}
\end{equation}
For finite support, the atomic formula is already rational and is
exactly recovered once the degrees are sufficient.
\end{proposition}

\begin{proof}
Use $M$-node Gaussian quadrature for $\mu$. Its nodes are the roots of
the monic orthogonal polynomial of degree $M$. Its orthogonality
equations use precisely moments through degree $2M-1$; infinite support
makes the moment matrix positive definite. All $M$ roots lie in the
interior of the convex hull of the support. Otherwise multiplication
by the product of the interior sign-changing factors would contradict
orthogonality by producing a nonzero integrand of one sign.

Define quadrature weights by integrating the Lagrange interpolation
polynomials at those nodes. Division by the orthogonal polynomial
proves exactness through degree $2M-1$. Applying exactness to the square
of a Lagrange polynomial proves positivity of its weight; applying it
to the constant gives total mass $m_0$. Expanding $(1+tx_j)^{-1}$
formally matches the first $2M$ coefficients. The combined numerator
and denominator have degrees at most $M-1$ and $M$, proving the Pad\'e
claim. This is the classical quadrature correspondence \cite{AllenChui}.

For the error, write
\[
 \frac1{1+tx}=\frac1{1+tR/2}
 \frac1{1+\dfrac{t(x-R/2)}{1+tR/2}}.
\]
The inner ratio has magnitude at most $tR/(2+tR)<1$. Its geometric
polynomial of degree $2M-1$ therefore approximates the integrand
uniformly in $x$ with error at most
\[
 \left(\frac{tR}{2+tR}\right)^{2M}.
\]
The geometric-remainder denominator cancels the leading factor since
$(1+tR/2)(1-tR/(2+tR))=1$. Integration and quadrature agree on the
polynomial and are both positive functionals of mass $m_0$. Their
remaining errors prove \eqref{eq:stieltjes-uniform}. Finite support
uses the exact atomic formula; zero mass is trivial.
\end{proof}

For $0<\varepsilon<2m_0$, a sufficient integer order is
\[
 M\ge\frac{\log(2m_0/\varepsilon)}
 {2\log\bigl((2+RT)/(RT)\bigr)}.
\]
There is no restriction $RT<1$. For example,
\[
 F(t)=\int_0^1\frac{dx}{1+tx}
 =\frac{\log(1+t)}{t},\qquad F(0)=1,
\]
has a Taylor series failing for $t>1$, while the error bound holds on
every finite $[0,T]$. For $1/(1+t)$ the $[0/1]$ Pad\'e approximant
is already exact, although its Taylor series also fails for $t>1$.

\subsection{Unbounded support, determinacy, and zero Taylor radius}
\label{app:determinate-stieltjes}

Suppose instead that $\mu$ is a finite positive measure on $[0,\infty)$
with every moment finite and no other positive measure on this half-line
having those moments. This is \emph{determinacy of the Stieltjes moment
problem}. The same $[M-1/M]$ Pad\'e approximants converge uniformly
to its actual Stieltjes integral on each $[0,T]$. The underlying theorem
appears in Simon \cite[Theorem 6]{SimonMoments}.

Here is a proof including $t=0$. The positive Gaussian quadrature
measure $\mu_M$ has mass $m_0$ and first moment $m_1$, so
\[
 \mu_M([A,\infty))\le m_1/A.
\]
Subsequences thus have weakly convergent further subsequences. For each
fixed $p\ge1$, quadrature eventually matches moments $p$ and $p+1$, and
\[
 \int_{x>A}x^p\,d\mu_M(x)\le m_{p+1}/A.
\]
Every moment passes to the weak limit. Determinacy identifies that
limit with $\mu$, so the whole sequence converges weakly. Integrating
the bounded continuous function $(1+tx)^{-1}$ gives pointwise convergence
for $t>0$; at zero, mass is exact. Finally, for $s,t\ge0$,
\[
 |R_M(t)-R_M(s)|\le m_1|t-s|,
 \qquad |F(t)-F(s)|\le m_1|t-s|.
\]
A finite time mesh upgrades convergence to uniform convergence on $[0,T]$.

A sufficient determinacy condition is the Stieltjes Carleman criterion
\[
 \sum_{p=1}^{\infty}m_p^{-1/(2p)}=\infty.
\]
For example $m_p\le C A^p(2p)!$ suffices, as does the stronger bound
$m_p\le C A^p p!$ \cite{SimonMoments,BorghiWeniger}.
These conditions are sufficient, not necessary.

The Euler example is
\[
 F(t)=\int_0^\infty\frac{e^{-x}}{1+tx}\,dx,
 \qquad F(t)\sim\sum_{p=0}^{\infty}(-1)^p p!\,t^p.
\]
Dominated differentiation gives right derivatives
$F^{(p)}(0)=(-1)^p(p!)^2$. Its Taylor series has zero convergence radius.
But $m_p=p!$ satisfies determinacy, so genuine direct Pad\'e converges
uniformly on every finite positive interval \cite{BorghiWeniger}.
Its ordinary Borel transform is particularly simple:
\[
 B(\xi)=\sum_{p\ge0}(-1)^p\xi^p=\frac1{1+\xi}
 \quad\hbox{near zero},\qquad
 F(t)=\int_0^\infty e^{-u}B(tu)\,du.
\]
The $[0/1]$ Pad\'e approximation of $B$ is exact. Its Laplace integral
is the nonrational function $F$. Thus direct Pad\'e and Borel--Pad\'e
both work here, but they are different constructions.

\subsection{Transferring uniform Borel-plane approximation to time}
\label{app:borel-transfer}

Fix $\sigma>0$ and suppose the \emph{actual} target satisfies
\begin{equation}
 F(t)=\int_0^\infty e^{-u}B(tu^\sigma)\,du,\qquad 0\le t\le T,
 \label{eq:borel-actual-F}
\end{equation}
where near zero
\[
 B(\xi)=\sum_{p\ge0}\frac{a_p}{\Gamma(1+\sigma p)}\xi^p.
\]
Identity \eqref{eq:borel-actual-F} is an assumption about the actual
function, not merely about its coefficients. Suppose rational
approximants $B_M$ from finite coefficients converge uniformly to $B$
on every finite positive-ray segment. A valid Pad\'e theorem for $B$
can supply this assumption.

Assume, for $C>0$ and $b\ge0$,
\[
 |B(\xi)|\le C e^{b\xi^{1/\sigma}},\qquad \xi\ge0.
\]
Fix $T$ with $\lambda=1-bT^{1/\sigma}>0$. For a cutoff $U$, define
\[
 F_{M,U}(t)=\int_0^U e^{-u}B_M(tu^\sigma)\,du.
\]
Splitting the exact integral at $U$ gives
\begin{equation}
 \sup_{0\le t\le T}|F_{M,U}(t)-F(t)|
 \le \sup_{0\le\xi\le TU^\sigma}|B_M(\xi)-B(\xi)|
       +\frac C\lambda e^{-\lambda U}.
 \label{eq:borel-cutoff}
\end{equation}
On the finite part all arguments stay in $[0,TU^\sigma]$ and the
weight has integral at most one. On the tail,
$e^{-u}|B(tu^\sigma)|\le Ce^{-\lambda u}$. Choose $U$ first, then $M$.
This supplies finite order for each requested accuracy on the fixed interval.

More generally it suffices that the true tail obey
\[
 \sup_{t\le T}\int_U^\infty e^{-u}|B(tu^\sigma)|\,du\longrightarrow0.
\]
No common growth bound on all $B_M$ is needed for the cutoff method.
If full Laplace integrals of $B_M$ are used instead, a sufficient
additional hypothesis is
\[
 |B_M(\xi)|+|B(\xi)|\le C e^{b\xi^{1/\sigma}}
\]
uniformly in $M$. The same split proves uniform convergence of the full
reconstructions. Compact convergence alone does not control these tails,
or justify an arbitrary growing cutoff $U=U_M$.

For a finite expression, partition $[0,U]$, choose a node $u_j$ in
each cell, and set $w_j=\int_{\text{cell }j}e^{-u}du$. Then
\[
 \widetilde F_{M,U,J}(t)=\sum_{j=1}^{J}w_jB_M(tu_j^\sigma).
\]
Continuity on the compact rectangle gives uniform quadrature convergence.
If $h$ is the largest cell length, its error is at most
\[
 T\sup_{0\le\xi\le TU^\sigma}|B_M'(\xi)|
 \begin{cases}
 h^\sigma,&0<\sigma\le1,\\
 \sigma U^{\sigma-1}h,&\sigma\ge1.
 \end{cases}
\]
For fixed $U,M$ this is finite. Choose tail, rational approximation and
quadrature errors each below $\varepsilon/3$. The result is a finite
rational function of $t$ with error below $\varepsilon$, computed from
initialization data and a declared quadrature rule. It need not be the
classical Pad\'e approximant of $F$: cutoff and quadrature can change the
exact coefficient matching. Qualitative convergence gives existence of
a sufficient order; a certified stopping rule additionally needs
quantitative bounds or a known convergence modulus.

\subsection{Stieltjes structure of the Borel transform}
\label{app:borel-stieltjes}

Suppose
\[
 B(\xi)=\int_{[0,R]}\frac{d\nu(x)}{1+\xi x},
 \qquad \nu\ge0,\qquad \nu_0=\nu([0,R])<\infty.
\]
Let $B_M=[M-1/M]_B$ and assume \eqref{eq:borel-actual-F}.
Positivity gives $0\le B,B_M\le\nu_0$ along the nonnegative ray.
Their full reconstructions
\[
 F_M(t)=\int_0^\infty e^{-u}B_M(tu^\sigma)\,du
\]
therefore satisfy, for every $U>0$ and finite $T$,
\begin{equation}
 \sup_{0\le t\le T}|F(t)-F_M(t)|
 \le 2\nu_0\left[
 \left(\frac{RTU^\sigma}{2+RTU^\sigma}\right)^{2M}+e^{-U}
 \right].
 \label{eq:borel-stieltjes-error}
\end{equation}
On $[0,U]$ apply \eqref{eq:stieltjes-uniform} to $B$; on the tail
use $|B-B_M|\le2\nu_0$. First choose $U$, then $M$. For example,
when $0<\varepsilon<4\nu_0$, take $U=\log(4\nu_0/\varepsilon)$ and
\[
 M\ge
 \frac{\log(4\nu_0/\varepsilon)}
 {2\log\bigl((2+RTU^\sigma)/(RTU^\sigma)\bigr)}.
\]
This gives error at most $\varepsilon$. Finite quadrature can supply
a rational expression. The original coefficients may grow like
$\Gamma(1+\sigma p)$, allowing zero Taylor radius.

There is also a direct Pad\'e conclusion when $0<\sigma\le2$.
Tonelli's theorem rewrites the actual function as
\[
 F(t)=\int_{[0,R]}\int_0^\infty
 \frac{e^{-u}}{1+tu^\sigma x}\,du\,d\nu(x)
 =\int_{[0,\infty)}\frac{d\rho(v)}{1+tv},
\]
where $\rho$ is the image of $d\nu(x)e^{-u}du$ under $v=u^\sigma x$.
Its moments satisfy
\[
 \int v^p\,d\rho(v)
 =\Gamma(1+\sigma p)\int x^p\,d\nu(x)
 \le\nu_0R^p\Gamma(1+\sigma p).
\]
Stirling's estimate bounds the Carleman terms below by a constant times
$p^{-\sigma/2}$. Their sum diverges for $0<\sigma\le2$, so $\rho$
is determinate. Appendix \ref{app:determinate-stieltjes} then gives
uniform convergence of the genuine $[M-1/M]$ Pad\'e approximants
of the original formal series on every $[0,T]$.
For $\sigma>2$ this particular determinacy argument stops; the
Borel--Pad\'e bound \eqref{eq:borel-stieltjes-error} remains valid.

\subsection{Conditions directly on the initialization coefficients}
\label{app:coefficient-tests}

For a proposed support bound $R>0$, put
\[
 s_p=\frac{(-1)^pa_p}{\Gamma(1+\sigma p)R^p}.
\]
The Hausdorff criterion says these are moments of a finite positive
measure on $[0,1]$ exactly when
\[
 \sum_{j=0}^{k}(-1)^j\binom{k}{j}s_{p+j}\ge0
 \qquad\hbox{for every }p,k\ge0.
\]
In one direction this expression is the integral of $x^p(1-x)^k$;
the converse is Hausdorff's theorem \cite{LiuPego}.
Rescaling by $x\mapsto Rx$ gives precisely the required compact
Stieltjes representation of $B$, including finite support and zero mass.

Another sufficient test from the earlier discussion is to put
$m_p=(-1)^pa_p/\Gamma(1+\sigma p)$ and require
\[
 (m_{i+j})_{i,j=0}^{q}
 \quad\hbox{and}\quad
 (m_{i+j+1})_{i,j=0}^{q}
\]
to be positive definite for every $q$, together with $m_p\le C R^p$.
The Stieltjes existence theorem gives a positive measure on
$[0,\infty)$ \cite{SimonMoments}. Any mass above $R$ would violate
the growth bound at high order, so its support lies in $[0,R]$.
The strict matrix criterion handles infinite support; finite support
can instead use its atomic formula or the Hausdorff criterion.

These are all-order conditions. Finitely many successful sign or
matrix tests are evidence, not a proof. Neither criterion identifies a
separately specified smooth observable with its Borel sum; that remains
the distinct identity \eqref{eq:borel-actual-F}.

\subsection{Other genuine Pad\'e guarantees and the limits of a general claim}
\label{app:pade-other}

Positivity is sufficient, not necessary. The de Montessus theorem
assumes $F$ is analytic near zero and meromorphic in $|t|<r$ with
exactly $q$ poles, counted with multiplicity. The fixed-denominator
row $[N/q]_F$ then converges uniformly on compact subsets avoiding
the poles as $N\to\infty$ \cite{Cacoq}. It gives the finite-order
guarantee on $[0,T]$ when $T<r$ and the interval avoids those poles.
Nearby complex poles can limit the Taylor radius while being captured
by the denominator. This is a row theorem, not an arbitrary-diagonal theorem.

For much wider classes with branch points, Stahl's theorem gives
convergence in logarithmic capacity in an appropriate continuation
domain \cite{Stahl}. That is weaker than maximum-error convergence.
Small exceptional sets may contain spurious poles, and one pole on
the target interval destroys a uniform bound. Analyticity or Borel
summability alone should not be presented as ensuring uniform
convergence of an arbitrary diagonal Pad\'e sequence.

The universal conjecture that a suitable diagonal subsequence always
converges locally uniformly throughout a meromorphy disk is false;
Lubinsky and Buslaev gave counterexamples \cite{Lubinsky,Buslaev}.
This does not say that every selectable approximant fails on each
particular real interval. It limits the blanket compact-uniform claim.

A transparent sufficient condition avoids positive measures. Suppose
$B$ is analytic on a connected complex neighborhood of zero and the
target interval. If the rational approximants are holomorphic there,
locally uniformly bounded, and match the coefficients at zero to
orders tending to infinity, then they converge locally uniformly to $B$.
Montel's theorem gives a convergent further subsequence from any
subsequence; Cauchy's formula passes derivatives to its limit; the
identity theorem identifies that limit with $B$. All subsequential
limits agree. The local boundedness is a substantive assumption,
not a consequence of coefficient matching. Related results upgrade
capacity convergence when poles are excluded from an open complex
neighborhood \cite{BeckermannMatos}; excluding poles only on the
real line is not the same condition.

\subsection{A constructive rational method from a known analytic domain}
\label{app:conformal-rational}

If a broader finite-jet rational method is acceptable, the continuation
can be proved without a generic diagonal Pad\'e theorem. Suppose $B$
is analytic in $\operatorname{Re}\xi>-a$, with $a>0$. The maps
\[
 w=\frac{\xi}{\xi+2a},\qquad \xi=\frac{2aw}{1-w}
\]
send the half-plane to the unit disk and back. Therefore
\[
 G(w)=B\!\left(\frac{2aw}{1-w}\right)=\sum_{p\ge0}g_pw^p
\]
is analytic for $|w|<1$. With $b_p=a_p/\Gamma(1+\sigma p)$,
\[
 g_0=b_0,\qquad
 g_p=\sum_{j=1}^{p}b_j(2a)^j\binom{p-1}{j-1}\quad(p\ge1).
\]
Thus the first $M+1$ coefficients use only the first $M+1$ Borel
coefficients. Define
\[
 R_M(\xi)=\sum_{p=0}^M g_p\left(\frac{\xi}{\xi+2a}\right)^p.
\]
Its only possible pole is $-2a$. On a fixed segment $[0,A]$, choose
$q=A/(A+2a)<\rho<1$. Cauchy's estimate gives
\[
 \sup_{0\le\xi\le A}|B(\xi)-R_M(\xi)|
 \le\max_{|w|=\rho}|G(w)|
 \frac{(q/\rho)^{M+1}}{1-q/\rho}.
\]
This tends to zero geometrically. A known bound on the maximum makes
the finite-order estimate effective. For example, for $C>0$ and $b\ge0$,
the stronger bound $|B(\xi)|\le C e^{b|\xi|}$ on the half-plane gives
\[
 \max_{|w|=\rho}|G(w)|
 \le C\exp\!\left(\frac{2ab\rho}{1-\rho}\right).
\]
For ordinary Borel summation ($\sigma=1$), that bound also controls
the true Laplace tail when $bT<1$.

These are Taylor truncations in a conformal coordinate. They match
$M+1$ coefficients, rather than the $2M+1$ conditions for conventional
$[M/M]$ Pad\'e, so that distinction must be preserved. The wider
conformal reconstruction framework is developed by Costin--Dunne
\cite[Theorem 8]{CostinDunne}. Combined with cutoff and quadrature,
this gives uniform approximation from finite initialization data under
the stated analytic-domain, tail, and identification assumptions.

\subsection{What these results would add to population gradient flow}

The main theorem supplies the actual population trajectory on a fixed
positive interval. Appendix \ref{app:borel-population} specifies the
additional reconstruction problem. For one input, the function to
reconstruct is $\kappa(f(0)+s)$: its independent argument is the
output increment $s$, not time. Uniform approximation transfers to
the scalar prediction equation. With several inputs, a separate
reconstruction of $K_{ab}(t)$ can be used with the weighted
prediction comparison in that appendix.

The approximation argument would be completed by proving the
observable's actual Borel identity and, for example, a Stieltjes
representation, an applicable meromorphic Pad\'e theorem, or a suitable
analytic domain with a declared rational construction. None of those
conditions is established here for the general nonlinear population
kernel. Stieltjes positivity is one sufficient route, not an assumption
of the original broader Borel conjecture. Once the reconstruction
premises hold on a fixed interval, the required finite order depends
on accuracy and that interval, not on width or the vanishing GD step.
